% 复旦大学本科毕业论文
% 基于 fduthesis 模板 (https://github.com/stone-zeng/fduthesis)

\documentclass[type=bachelor,oneside]{fduthesis}

\fdusetup{
  style = {
    cjk-font = windows,
    bib-backend = bibtex,
    bib-style = numerical,
    bib-resource = {refs.bib},
    declaration-page = {declaration.pdf},
  },
  info = {
    title = {基于小样本学习的机器人决策规划技术研究},
    title* = {Research on Robot Decision Planning Technology Based on Few-Shot Learning},
    author = {任立德},
    supervisor = {戈维峰\quad 副教授},
    major = {软件工程},
    department = {计算与智能创新学院},
    student-id = {22302010001},
    keywords = {三维场景理解；数据合成；视觉问答；游戏逻辑驱动；仿真验证},
    keywords* = {3D scene understanding; data synthesis; visual question answering; game-logic-driven; simulation verification},
    clc = {TP391},
  }
}

\usepackage{enumitem}
\usepackage{booktabs}
\usepackage{longtable}
\usepackage{float}
\usepackage{tikz}
\usetikzlibrary{arrows.meta, positioning, decorations.pathreplacing}

% 按本科论文撰写规范覆盖标题字号
% fduthesis 使用 \sffamily 对应黑体, 不使用 \heiti
% 章标题: 小二号(18pt)黑体居中
\ctexset{chapter={
    format      = {\centering\sffamily\zihao{-2}},
    beforeskip  = {20pt},
    afterskip   = {18pt},
}}
% 一级标题(1.1): 四号(14pt)黑体
\ctexset{section={
    format      = {\sffamily\zihao{4}},
    beforeskip  = {12pt},
    afterskip   = {6pt},
}}
% 二级标题(1.1.1): 小四(12pt)黑体
\ctexset{subsection={
    format      = {\sffamily\zihao{-4}},
    beforeskip  = {6pt},
    afterskip   = {3pt},
}}

\begin{document}

\frontmatter
\tableofcontents

\begin{abstract}
三维场景理解任务面临高质量多模态标注数据稀缺的核心困境。人工标注成本高昂——单个室内场景（约 70 个物体）的完整标注可能需要数小时的专家劳动——导致标注数据的规模和多样性严重受限。现有自动生成方法虽然降低了成本，但在关系推理精度和跨场景泛化性上存在明显不足：视觉语言模型倾向于生成光照、氛围等场景级感官描述，而非精确的物体类别和空间关系信息。

针对上述问题，本文提出了一种融合游戏逻辑驱动合成、真实场景理解和仿真真值验证的三层数据自动标注方法。首先，将 Code2Logic 的模板化数据合成范式从二维游戏领域迁移至三维场景，以 Concept-Graphs 为框架整合 GradSLAM、Grounded-SAM、CF-SLAM、LLaVA v1.5 和 GPU 加速的大语言模型推理，构建了从 RGB-D 图像序列到结构化 QA 数据集的端到端六步自动处理管线。通过专门设计的结构化 prompt，将物体类别识别率从接近 0 提升至约 85\%。其次，针对 Ollama 推理引擎存在的 CUDA 版本不兼容问题，通过 llama-cpp-python 源码编译方案将推理吞吐量从 3 tok/s 提升至 99 tok/s（33 倍提升），单场景 LLM 处理时间从 15 分钟降至 50 秒。最后，集成 NVIDIA Isaac Sim 仿真环境构建了合成真值验证系统，通过程序化场景生成和几何规则计算，实现了标注精度的可量化评估。

实验结果表明，本文方法在三类数据源（30 款 2D 游戏、7 个真实室内场景、20 个合成仿真场景）上累计产出 2395 条结构化 QA 数据，覆盖 Target Perception、Spatial Reasoning、Scene Comprehension、State Prediction 和 Strategy Optimization 共 5 种推理类型。跨源对比评估揭示了合成真值对推断标注的可量化校准能力——Isaac Sim 基于几何规则的空间关系准确率为 100\%，而 Concept-Graphs 依赖模型推断的准确率约为 70\%。5 款国内开源大语言模型的基准测试表明，Qwen2.5-3B 在速度和质量的平衡上为最优选择（86 tok/s，100\% 格式合法率）。

本研究为三维场景的自动化标注提供了一套低成本、可扩展、精度可验证的完整解决方案，对视觉语言模型训练数据的构建和多模态场景理解系统的开发具有一定的理论价值和工程参考意义。
\end{abstract}

\begin{abstract*}
Three-dimensional scene understanding faces the fundamental challenge of scarce high-quality multimodal annotation data. Manual annotation is prohibitively expensive—a single indoor scene with approximately 70 objects can require hours of expert labor—severely limiting the scale and diversity of training data for vision-language models. Existing automatic methods reduce cost but exhibit significant shortcomings in relational reasoning accuracy and cross-scene generalization. Specifically, current VLMs tend to generate scene-level sensory descriptions (e.g., lighting conditions, atmospheric qualities) rather than precise object-level semantic information.

To address these challenges, this thesis proposes a three-tier data annotation framework that integrates game-logic-driven synthesis, real-world scene understanding, and simulation-based ground-truth verification. The framework begins by migrating the template-based data synthesis paradigm of Code2Logic from 2D games to 3D scenes. Using Concept-Graphs as the integration framework, we construct an end-to-end six-step pipeline that combines GradSLAM (3D reconstruction), Grounded-SAM (2D segmentation), CF-SLAM (object-level mapping), LLaVA v1.5 (multimodal description), and GPU-accelerated LLM inference (description refinement and relation reasoning). A specially designed structured prompt increases the object category recognition rate from near zero to approximately 85\%. Furthermore, we identify and resolve a CUDA version incompatibility in the Ollama inference engine by source-compiling llama-cpp-python, achieving a 33$\times$ throughput improvement (3 $\rightarrow$ 99 tok/s) and reducing per-scene LLM processing time from 15 minutes to 50 seconds. Finally, we integrate NVIDIA Isaac Sim to construct a synthetic ground-truth verification system, enabling quantifiable evaluation of annotation accuracy through procedural scene generation and geometric rule computation.

Experimental results demonstrate that our method produces 2,395 structured QA pairs across three data sources (30 2D games, 7 real indoor scenes, and 20 synthetic simulation scenes), covering 5 reasoning types: Target Perception, Spatial Reasoning, Scene Comprehension, State Prediction, and Strategy Optimization. Cross-source comparative evaluation reveals the quantifiable calibration capability of synthetic ground truth—Isaac Sim achieves 100\% spatial relation accuracy via geometric rules, while Concept-Graphs LLM-inferred accuracy is approximately 70\%. Benchmarking of five Chinese open-source LLMs identifies Qwen2.5-3B as the optimal choice (86 tok/s, 100\% JSON validity).

This work provides a low-cost, scalable, and verifiable solution for automatic 3D scene annotation, offering both theoretical insights and practical engineering value for vision-language model training data construction and multimodal scene understanding system development.
\end{abstract*}

\mainmatter

\chapter{引言}
% 第1章 引言

\section{研究背景}

让机器理解三维室内场景——识别房间里的桌子、椅子、沙发，搞清楚谁在谁旁边、谁在谁上面——是计算机视觉和机器人领域的一个基础问题。这类任务需要标注数据：每张图中每个物体属于什么类别，它们之间的空间关系是怎样的。

人工标注三维场景的成本很高。一个包含约 70 个物体的中等规模房间，从头到尾标完需要几个小时的专家劳动，而且标注者的空间判断还会引入主观误差。视觉语言模型（VLM）可以自动生成描述，算是个省力的替代方案，但生成质量经常不过关。这类模型倾向于输出诸如"自然光洒满整个空间"这样的全局氛围描述，而不是"靠墙放着一张棕色木桌"这种具体的物体描述。Song 等人\cite{llm_planner} 的测试提供了一个量化参考：在 Stanford VLN-CE 长程导航任务中，当步骤超过 5 步时 VLM 的决策准确率下降 37.2\%，环境扰动让动作成功率再降 42.8\%。

Tong 等人\cite{code2logic} 从游戏领域找到了一条有意思的解决路径。他们提出的 Code2Logic 方法基于一个简单的观察：游戏源代码里的坐标数组、碰撞检测逻辑、胜负判定规则本身就是天然的结构化标注。让大语言模型读游戏代码，按模板生成问答对，就能零边际成本地无限产出完美的标注数据。他们构建的 GameQA 数据集包含了 30 款游戏和 14 万条问题，训练后模型在 7 个视觉基准上平均提升了 2.33 个百分点。

这就引出了一个自然的迁移问题：能不能把游戏代码驱动的这套自动化方法搬进真实三维场景？

这个迁移至少面临三个困难。游戏逻辑是显式的，物体在哪、能做什么都由代码精确指定；真实场景的逻辑是隐式的，物体要从 RGB-D 数据里分割和识别。游戏画面是完美渲染的，没有噪声；真实的传感器数据有噪声、遮挡和光照变化。游戏里的每一步操作都可以和代码对比来验对错，真实场景没有这种天然的正确答案。本文的工作围绕这三个困难展开：把 Code2Logic 的模板化方法从 2D 游戏迁移到 3D 场景，构建一条端到端的自动标注管线，再用合成真值系统来验证标注精度。

\section{本文工作}

本文搭建了一条三层数据合成链路。底层是 Code2Logic，30 款游戏的 1397 条 QA 证明了从代码逻辑到模板再到数据这条路可行。中层是 Concept-Graphs，整合了 GradSLAM、Grounded-SAM、CF-SLAM、LLaVA 和 GPU 推理，串成了一条从 RGB-D 图像序列到结构化 QA 数据集的六步流水线。上层是 Isaac Sim，生成 20 个合成场景，每个物体的坐标、类别和尺寸都是精确已知的，用作验证标注精度的金标准。三层递进：方法在底层验证，在中层迁移到真实场景，在上层被合成真值校准。

本文修了一条端到端的自动标注管线。中间六个步骤依次是：GradSLAM 做 3D 重建，Grounded-SAM 做物体分割，CF-SLAM 做三维建图，LLaVA v1.5 为每个物体裁剪区域生成自然语言描述，llama.cpp 在 GPU 上批量优化描述并推理物体间的空间关系，最后按模板生成 QA 数据集。管线的关键改进在第五步的 prompt 设计。LLaVA 的原始描述有两个常见问题：太笼统（"自然光洒满空间"）和输出被截断（"es are all white..."）。重写后的 prompt 要求模型显式输出物体类别（door、table、sofa 等），同时明确禁止了 light、atmosphere、depth 等十几个场景级描述词。这个改动把物体类别识别率从接近零提到了约 85\%。

本文还对 GPU 推理做了一次深度优化。Ollama 自带的 CUDA 12.8 库与服务器驱动 565.57 不兼容，导致推理全部落在 CPU 上，实测吞吐量只有约 3 tok/s。换用源码编译的 llama-cpp-python（链接系统 CUDA 12.6 库），qwen2.5-3B 模型的全部 29 层加载到单张 RTX 4090 显存上。优化后吞吐量达到 99 tok/s，单个场景的 LLM 处理时间从 15 分钟降到了 50 秒。在 5 款国产开源模型上的横向对比表明，Qwen2.5-3B 在速度（86 tok/s）、质量（100\%）和格式合法性（100\%）三个维度上综合最优，显存仅需 1.9 GB。

本文建了一套跨源评估体系。从数据量、推理类型分布、关系类型分布、标注精度和管线效率五个维度，对比了 Code2Logic（游戏合成，1397 QA）、Concept-Graphs（真实 3D 推断，370 QA）和 Isaac Sim（合成 3D 真值，795 QA）三源数据。三类数据累积 2395 条 QA，覆盖五种推理类型。彼此之间是互补关系，不是替代关系：Code2Logic 提供涉及多步规划的 Strategy Optimization 类 QA（这是静态三维场景里不出现的推理类型），Isaac Sim 贡献了大量高精度空间关系（610 条，准确率 100\%），Concept-Graphs 则用开放词汇模型弥补了 Isaac Sim 只能识别预设家具类别的局限。

\section{关键术语与问题边界}

本文中的"标注数据"特指结构化的视觉问答格式，每条包含一个关于场景中物体的问题、一个答案和一个类型标签。选择 VQA 格式是因为它目前是视觉语言模型训练和评估的主流范式，产出的数据可以直接用于下游微调，而且问题-答案的结构天然适合模板化生成。管线的自动化程度是从 RGB-D 数据到 QA 数据集的自动转化，中间不需要人工标注或人工设计规则，但管线依赖公开的 RGB-D 数据集作为输入，大语言模型的 prompt 模板也需要人工设计。标注质量的评测目前依赖 Isaac Sim 合成场景的 ground truth 作为推断标注的验证基准，加上作者对部分 QA 的人工定性抽查。大规模人工评估和下游 VLM 训练验证尚未完成，这两个局限在第五章中有具体讨论。本文方法的设计和验证均限于室内静态场景，室外场景、工业环境和包含动态物体的场景不在当前实验范围内。

\section{论文结构}

第二章梳理相关技术：VLM、CoT 推理、Code2Logic、三维场景图构建和合成数据方法。第三章展开方法细节，包括三层架构、六步管线、prompt 设计、GPU 加速方案和 Isaac Sim 验证系统。第四章报告实验：7 个场景的标注产出、描述质量前后对比、prompt 消融实验、5 模型选型基准、管线效率分析和三源数据多维对比。第五章总结全文，讨论局限性和后续方向。


\chapter{相关工作}
% 第2章 相关工作

\section{视觉语言模型}

视觉语言模型（Vision-Language Models，VLMs）是连接计算机视觉与自然语言处理的桥梁技术，通过将视觉编码器与大语言模型对齐，实现跨模态图像理解和文本生成能力。近年来，该类模型的快速发展为三维场景的自动化描述提供了技术基础。

在开源视觉语言模型方面，LLaVA（Large Language and Vision Assistant）\cite{llava} 是代表性的工作之一。其核心架构由 Vision Transformer（ViT）视觉编码器和 LLaMA 语言解码器组成，通过一个简单的线性投影层将视觉特征映射至语言空间。LLaVA v1.5 7B 版本在视觉指令微调后，能够对任意图像区域生成自然语言描述。然而，LLaVA 的描述质量高度依赖于输入图像的质量和分辨率。在本文的三维物体裁剪场景中，由于物体裁剪区域往往较小（如仅占原始图像 5\%—15\%），且包含大量背景噪声，LLaVA 生成的描述常常偏向于整体场景氛围而非物体本身属性，这一局限性成为本文后续 prompt 优化的直接动机。

Qwen2.5 系列模型\cite{qwen2_5} 是阿里巴巴 Qwen 团队发布的开源大语言模型家族，涵盖 0.5B 至 72B 参数的多个规模变体。其中 Qwen2.5-3B 版本在推理能力和资源消耗之间取得了优秀的平衡：在 MMLU 基准上达到 65.6 分，同时推理速度远快于 7B 及以上版本。本文在推理引擎选型实验中系统测试了 Qwen2.5 系列的三个变体（3B、7B、14B），验证了小模型在结构化数据合成任务中的充分性。Qwen2.5 的 GGUF 量化格式（Q4\_K\_M）进一步将模型压缩至 1.9—4.4 GB，使得在单张消费级 GPU 上的本地化部署成为可能。

An 等人（ICLR 2025）的研究表明，将文本等额外模态信息与三维点云数据融合，可以显著提升 few-shot 三维语义分割的性能\cite{multimodal_3d}。该工作使用预训练的 CLIP 模型提取文本特征，通过跨模态注意力机制将语言先验知识注入点云分割网络，在 ScanNet 和 S3DIS 数据集上均取得了 state-of-the-art 的 few-shot 性能。这一发现提示了多模态融合对三维理解任务的重要性，也为本文后续引入多模态描述融合提供了理论基础。

\section{思维链推理}

思维链（Chain-of-Thought，CoT）提示是一种通过显式引导语言模型生成中间推理步骤来提升复杂任务性能的方法。其核心思想是模仿人类解决问题的过程——将复杂问题拆解为多个可处理的子步骤，逐步推导出最终答案。在 GSM8K 数学数据集上，PaLM-540B 使用 8 个 CoT 示例后准确率从 34\% 提升至 57\%，超越了当时微调后的 GPT-3 175B，证明了 CoT 在多步推理任务中的有效性。

然而，传统 Few-shot CoT 存在明显局限。首先，人工设计的 CoT 示例不仅耗时耗力，且示例质量直接影响推理效果。Zhang 等人提出的 Auto-CoT\cite{auto_cot} 通过聚类采样自动生成演示：将待处理问题聚类为若干类别，从每个类别中选择代表性样本，使用语言模型为这些样本自动生成推理链。Auto-CoT 在多个推理基准上接近甚至超越了人工设计 CoT 的效果，验证了自动化示例生成的可行性。

Cheng 等人对新一代强大语言模型上的 CoT 提示进行了系统性重新审视\cite{chain_of_thought}。该研究发现了一个重要的反直觉现象：对于 Qwen2.5 等强大语言模型，Zero-Shot CoT 在 GSM8K 和 MATH 数据集上的推理准确率普遍高于传统 Few-shot CoT。进一步的替换实验揭示：传统 CoT 示例的核心作用仅是规范输出格式（如引导模型使用特定格式包裹答案），而非传递推理知识。注意力机制可视化进一步证实：模型在解码时几乎不关注示例区域，而是聚焦于问题指令本身。这一发现为本文采用零样本结构化 prompt 进行场景描述优化提供了直接的理论支撑——对于能力足够强大的语言模型，精心设计的格式约束比提供示例更有效。

Zhou 等人探索了使用大语言模型为软件架构决策生成设计理由的可行性\cite{design_rationale}。其研究结果表明零样本方法的平均精确度（0.278）和 F1 值（0.381）均优于 CoT 方法（0.237 和 0.351），且推理速度最快、资源消耗最低。CoT 生成的论证虽然数量增加 26.7\%，但其中许多论证流于表面、缺乏深度分析，进一步验证了零样本推理在结构化输出任务中的优势。

\section{Code2Logic：游戏逻辑驱动的数据合成}

Code2Logic 由 Tong 等人提出（NeurIPS 2025），是本文方法论的核心来源\cite{code2logic}。该方法创造性地将游戏源代码作为高质量推理数据的自动化生成源，提出了三阶段数据合成框架。

第一阶段为游戏代码构建。使用大语言模型根据游戏规则描述生成可执行的游戏源代码。以推箱子为例，仅需单行自然语言提示，LLM 即可输出包含状态空间定义、合法动作集合、胜负判定条件的完整 Python 代码，建立了从自然语言规则到程序语义的精确映射。

第二阶段为问答模板设计。从游戏逻辑中提取三类渐进式推理任务：（1）Target Perception（目标感知）——聚焦静态元素识别；（2）State Prediction（状态预测）——模拟动作执行后的状态变化；（3）Strategy Optimization（策略优化）——在多步推理中寻求最优解。这三类任务构成认知梯度的渐进式推理挑战。

第三阶段为数据引擎构建。将游戏核心代码作为数据生成引擎，通过自动化参数采样和状态快照，批量生成带视觉轨迹的问答对。单款游戏平均 7.5 小时的初始开发投入后，数据引擎可零边际成本无限生成新样本，单位标注成本仅为人工标注的 1/50。

Code2Logic 构建的 GameQA 数据集涵盖 30 种游戏、158 项认知任务和 14 万问题。实验揭示了几项关键洞察：（1）仅用 GameQA 训练的 VLM 在 7 个主流视觉基准上平均提升 2.33\%；（2）使用 20 款游戏训练的模型比 4 款游戏的域外任务准确率提升 1.8\%，验证了游戏多样性对泛化能力的驱动作用；（3）GRPO 强化学习策略使视觉感知与文本推理错误率分别降低 19.2\% 和 11.5\%。

Code2Logic 的核心局限在于其完全依赖游戏代码的显式逻辑结构——游戏世界中物体的位置和交互规则均由代码明确定义。而真实物理场景中，这些信息需要通过传感器数据和视觉推理间接获取，噪声、遮挡和光照变化都会引入不确定性。本文的工作正是试图桥接这一差距：将模板化数据合成方法论从二维游戏世界迁移至三维真实场景。

\section{三维场景图构建}

三维场景图（3D Scene Graph）将场景表示为结构化图——节点代表物体实例，边代表空间或语义关系。与二维场景图不同，三维场景图需要处理深度信息、几何约束和视角变化等额外挑战。

Concept-Graphs\cite{conceptgraphs} 是三维场景理解领域的重要工作，通过整合多种感知模块实现从 RGB-D 到三维场景图的端到端构建。其三个核心组件分别为：GradSLAM\cite{gradslam} 提供可微分的 SLAM 前端，从 RGB-D 序列和相机轨迹重建三维场景几何；Grounded-SAM\cite{grounded_sam} 组合 Grounding DINO 的开放词汇检测和 SAM 的精细分割，实现任意物体的检测与分割；CF-SLAM 将二维分割结果通过相机参数投影至三维空间，经重叠检测、相似度匹配和 DBSCAN 聚类，构建物体级场景图。

Replica 数据集\cite{replica} 由 Straub 等人发布，包含 18 个高精度三维重建的室内场景。每个场景提供约 2000 帧 RGB 图像（640×480）、对应深度图和由运动捕捉系统采集的相机轨迹。本文使用 room0、room1 和 office0—4 共 7 个场景，覆盖起居室和办公室两类室内环境。

\section{合成数据与域随机化}

合成数据的核心思想是通过计算机图形学自动生成带标注的训练数据。其在计算机视觉领域的核心优势包括：标注自动且完美准确、场景参数完全可控、可无限扩展。

域随机化（Domain Randomization）在合成数据基础上进一步提升模型泛化能力：在仿真环境中对非本质属性（光照、颜色、纹理等）进行随机扰动，使模型学到对这些变化不敏感的鲁棒特征。本文在 Isaac Sim 中应用了三种域随机化：物体位置随机化、材质颜色随机化、光照参数随机化。

NVIDIA Isaac Sim\cite{isaac_sim} 是面向机器人仿真和数据合成的综合性平台，基于 USD（Universal Scene Description）格式，支持 RTX 光线追踪渲染和无头模式 Vulkan 后端。在本文中，Isaac Sim 被用作标注质量的验证工具——以合成场景中的精确 ground truth 为基准，对 Concept-Graphs 管线的推断结果进行可量化的精度评估。

\section{本章小结}

综上所述，视觉语言模型为三维场景理解提供了基础能力，思维链推理研究明确了零样本结构化输出的优越性，Code2Logic 验证了模板化数据合成的工程可行性，三维场景图技术提供了从传感器数据到结构化表示的完整链路，合成数据方法则为标注质量的闭环验证奠定了基础。本文的创新正在于将这些技术线交汇融合——将游戏逻辑驱动的合成范式从二维游戏世界迁移至三维真实场景，构建可验证的自动标注系统。


\chapter{方法}
% 第3章 方法

\section{整体架构}

本文的标注系统分三层，如图\ref{fig:architecture}所示。

% 图1: 三层系统架构
\begin{figure}[H]
\centering
\begin{tikzpicture}[node distance=0.8cm,
    box/.style={rectangle, draw=black, thick, minimum width=10cm, minimum height=1.2cm, align=center, font=\small}]
\node[box, fill=blue!8] (cl) {{\bf 数据合成层 -- Code2Logic}\\{游戏代码解析 $\to$ 模板填充 $\to$ 结构化 QA}};
\node[box, fill=green!8, below=of cl] (cg) {{\bf 场景理解层 -- Concept-Graphs}\\{SLAM $\to$ SAM $\to$ CF-SLAM $\to$ LLaVA $\to$ LLM $\to$ QA}};
\node[box, fill=orange!8, below=of cg] (is) {{\bf 真值验证层 -- Isaac Sim}\\{USD 场景生成 $\to$ 域随机化 $\to$ 几何真值导出}};
\node[box, fill=gray!10, below=of is] (out) {{\bf 输出: 结构化 QA 数据集 (2395 条, 5 种推理类型)}};
\draw[->, >=stealth, thick] (cl) -- (cg) node[midway, right] {\footnotesize 方法迁移};
\draw[->, >=stealth, thick] (cg) -- (is) node[midway, right] {\footnotesize 真值校准};
\draw[->, >=stealth, thick] (is) -- (out);
\end{tikzpicture}
\caption{系统三层架构。}
\label{fig:architecture}
\end{figure}


第一层是数据合成层。Code2Logic 从游戏源代码里提取逻辑结构，按模板生成 QA。这层的作用是证明方法论可行——不需要人工标注，从代码就能自动产出高质量问答数据。

第二层是场景理解层。把 Code2Logic 的模板化方法搬到真实 3D 场景。这里要处理的东西比游戏世界复杂得多——物体不是代码定义的坐标，而是从 RGB-D 数据里分割出来的；空间关系不是游戏规则写死的，要靠模型推断。

第三层是真值验证层。用 Isaac Sim 生成合成场景，每个物体的坐标、类别、尺寸都是精确已知的。以这些"标准答案"为基准，可以定量地算 Concept-Graphs 管线推断出来的标注到底有多准。

三层是递进关系：游戏方法在第 1 层被验证，搬到第 2 层解决真实问题，在第 3 层被校准。

\section{Concept-Graphs 标注管线}

\subsection{六步管线}

管线的数据流如图\ref{fig:pipeline}所示。

% 图2: 六步标注管线
\begin{figure}[H]
\centering
\begin{tikzpicture}[node distance=0.4cm,
    sbox/.style={rectangle, draw=black, thick, rounded corners=3pt, minimum width=1.8cm, minimum height=1.3cm, align=center, font=\footnotesize, fill=blue!5}]
\node[sbox] (s1) {RGB-D\\序列};
\node[sbox, right=of s1] (s2) {GradSLAM\\3D重建};
\node[sbox, right=of s2] (s3) {Grounded-\\SAM分割};
\node[sbox, right=of s3] (s4) {CF-SLAM\\3D建图};
\node[sbox, right=of s4] (s5) {LLaVA v1.5\\物体描述};
\node[sbox, right=of s5, fill=green!10] (s6) {llama.cpp\\GPU优化};
\node[sbox, right=of s6, fill=gray!15] (s7) {QA数据集};
\draw[->, >=stealth, thick] (s1) -- (s2);
\draw[->, >=stealth, thick] (s2) -- (s3);
\draw[->, >=stealth, thick] (s3) -- (s4);
\draw[->, >=stealth, thick] (s4) -- (s5);
\draw[->, >=stealth, thick] (s5) -- (s6);
\draw[->, >=stealth, thick] (s6) -- (s7);
\end{tikzpicture}
\caption{六步自动标注管线。}
\label{fig:pipeline}
\end{figure}


第一步是三维重建，用 GradSLAM\cite{gradslam} 从 RGB-D 序列和相机轨迹建点云地图，以 5 帧为步长处理约 400 个关键帧。第二步是二维分割，用 Grounded-SAM\cite{grounded_sam} 在每个关键帧上做检测和分割，这是最耗时的一步，单场景约 25 分钟。第三步是三维建图，用 CF-SLAM 把二维分割投影到三维空间，经重叠检测和 DBSCAN 聚类后输出每个物体的三维位置、包围盒和语义特征。

第四步是物体描述，用 LLaVA v1.5 7B\cite{llava} 对每个物体的裁剪区域生成自然语言描述。第五步是优化和推理，用 llama.cpp 在 GPU 上跑 qwen2.5:3b\cite{qwen2_5}，同时做两件事：优化 LLaVA 的原始描述（让它更像"物体标注"而非"场景描述"），推断物体之间的空间关系。第六步是按模板生成 QA 数据集。

\subsection{Prompt 设计}

LLaVA 输出的典型问题是两个：描述太笼统（"自然光洒满空间"），以及输出被截断（"es are all white..."）。针对这两个问题，本文设计了一套新的 prompt，核心改动有三条。

第一条是强制输出 class 字段——要求模型显式给出物体类别，例如 door、table、sofa。这直接解决了旧 prompt 输出"全是氛围、没有物体"的问题。第二条是设置禁用词汇表——light、atmosphere、depth、shadow 等十几个词被直接禁止。这是一个简单粗暴但极其有效的手段：模型生成这些词的概率被压到零，它就只能往"描述物体"的方向走。第三条是容错处理——输入太模糊时标记为 unknown，输入被截断时直接跳过。这避免了低质量数据流入下游。

Cheng 等人\cite{chain_of_thought} 的结论为这个设计提供了背书：对大模型来说，说清楚你要什么格式，比给一堆例子重要得多。

\subsection{GPU 加速}

原始管线用 Ollama 做本地推理，但 Ollama 内置的 CUDA 12.8 库（libggml-cuda.so）和服务器驱动（565.57.01，CUDA 12.7）不兼容。虽然 Ollama 声称把模型层"offload 到了 GPU"，但 nvidia-smi 上只显示 4 MiB 显存，实际计算全在 CPU 上——吞吐量只有约 3 tok/s，单场景处理要 15 分钟。

解决办法是换 llama-cpp-python。服务器上本来就装了 CUDA 12.6 完整工具链（nvcc 12.6、cmake 3.22、gcc 11.4），直接用这些系统库源码编译 v0.2.90。关键编译参数是 \texttt{-DGGML\_CUDA=on}，CUDA 架构设为 sm\_89。编译完成后，qwen2.5:3b GGUF 模型（Q4\_K\_M，1.9 GB）的 29 层全部加载到单张 RTX 4090 显存，开一个 OpenAI 兼容 API（端口 8080），管线里其它代码一行不用改。

编译中踩了三个坑。第一个是系统默认 nvcc 是 11.5，不支持 RTX 4090 的 sm\_89。通过 \texttt{-DCMAKE\_CUDA\_COMPILER} 指定 CUDA 12.6 的 nvcc 路径解决。第二个是缺少 OpenMP，用 \texttt{-DGGML\_OPENMP=off} 关掉。第三个是 CUDA 架构字符串格式——CMake 4.3 需要写成 \texttt{89} 而非 \texttt{sm\_89}。

优化后吞吐量 99 tok/s，显存占用 4.2 GB（7B 模型）或 1.9 GB（3B 模型）。

\section{Isaac Sim 验证系统}

\subsection{场景生成}

Isaac Sim 4.5 通过 pip 安装在服务器上，无头模式运行，Vulkan 渲染。场景生成完全程序化，不依赖图形界面。

设计了四类房型模板：起居室（6m×5m，10 个物体）、卧室（4.5m×4m，8 个物体）、厨房（5m×4m，8 个物体）、书房（4m×3.5m，7 个物体）。每类生成 5 个随机化变体，共 20 个场景。随机化覆盖三个维度：物体的位置在房间边界内随机扰动（碰撞检测保证间距 ≥0.5m），颜色从 3 种预设变体中随机选取，光照色温和强度连续随机采样。

\subsection{真值导出}

Isaac Sim 的 USD 场景图里，每个物体的坐标、类别和尺寸都是明确存储的，API 直接查询就能拿到标准答案，位置精度 0.01 米。空间关系用几何规则算：水平距离小于 1.5 米判定为 next to，1.5--3.0 米为 near，超出 3 米按方位角判断前后左右。因为计算过程是纯确定性的，准确率 100\%——不存在模型推断的不确定性。

场景图的导出格式和 Concept-Graphs 的标注模板一模一样，所以两边的数据可以直接对比。

\section{评估框架}

跨源评估从五个维度对比三类数据：（1）数据量——各源的 QA 总数和来源数；（2）QA 类型——Target Perception、Spatial Reasoning 等的分布；（3）关系类型——next to、above、behind 等的频次；（4）标注精度——类别准确率和关系准确率；（5）管线效率——数据获取方式和单场景耗时。


\chapter{实验}
% 第4章 实验

\section{实验设置}

\subsection{硬件环境}

实验在一台配备 7 张 NVIDIA GeForce RTX 4090（24 GB 显存/卡）、503 GB 系统内存的 Ubuntu 22.04 服务器上进行。NVIDIA 驱动版本为 565.57.01，CUDA 版本 12.7。大语言模型推理在单张 RTX 4090 上运行（通过 \texttt{CUDA\_VISIBLE\_DEVICES=0} 指定），LLaVA 视觉编码使用另一张独立 GPU。Isaac Sim 4.5 以无头模式运行，使用 Vulkan 渲染后端自动检测可用的 GPU。

\subsection{软件环境}

操作系统为 Ubuntu 22.04（glibc 2.35），Python 版本 3.10。核心依赖包括：PyTorch 2.5.1+cu124，transformers 4.37.2（LLaVA v1.5 兼容要求），llama-cpp-python v0.2.90（源码编译，链接 CUDA 12.6），NVIDIA Isaac Sim 4.5.0（pip 安装）。大语言模型推理使用 qwen2.5:3b GGUF 模型（Q4\_K\_M 量化），通过 llama-cpp-python 的 OpenAI 兼容 API 服务器（端口 8080）提供服务。

LLaVA v1.5 7B 的加载需要 transformers 4.37.2 的特定版本约束，与当前最新的 4.49.0 存在 API 不兼容（\texttt{forward()} 方法新增的 \texttt{cache\_position} 参数不被 LLaVA 的 \texttt{LlavaLlamaForCausalLM} 接受）。pip 依赖解析器不会自动处理这一不兼容，需要在环境搭建时手动指定版本。

\subsection{数据集}

本实验使用三类数据源：Replica 数据集\cite{replica} 的 7 个室内场景（room0、room1、office0—4），每个场景含约 2000 帧 RGB-D 图像（640 $\times$ 480）及相机轨迹；Isaac Sim 合成数据集包含 4 类房型各 5 个随机化场景变体（共 20 个场景），每场景含精确的物体坐标和类别真值；Code2Logic GameQA 数据集\cite{code2logic} 包含 30 款游戏的 1397 条 QA，覆盖 Target Perception（660 条）、State Prediction（470 条）和 Strategy Optimization（267 条）三种推理类型。

\section{游戏数据合成实验}

Code2Logic 数据集作为本文方法论的来源和对比基线，不涉及训练过程。1397 条 QA 覆盖 30 款不同类型的游戏，包括棋类（PyramidChess、tictactoe、ultra\_tictactoe）、解谜类（sokoban、sudoku、tangram）、动作类（pacman、space\_invaders、zuma）和策略类（minesweeper、klondike、rubiks\_cube）等。其中 rubiks\_cube 以 498 条 QA 贡献了单款游戏的最大数据量，minesweeper 以 180 条紧随其后。数据生成的平均速度约为 205 QA/s，单款游戏的初始开发投入约 7.5 小时后可无限生成新样本。

从推理类型的分布来看，Target Perception 占 47\%（660/1397），侧重于物体和元素的静态识别；State Prediction 占 34\%（470/1397），考察动作对游戏状态的影响预测；Strategy Optimization 占 19\%（267/1397），要求多步规划能力。后两类 QA（State Prediction 和 Strategy Optimization）涉及时序推理和规划，是静态三维场景数据中天然缺失的推理类型，这也提示了未来将动态场景引入标注管线的潜在价值。

\section{真实场景标注实验}

\subsection{跨场景标注结果}

使用优化后的 prompt 和 GPU 加速管线，在 7 个 Replica 场景上运行完整的六步标注流程。表 \ref{tab:scene_results} 汇总了各场景的详细标注产出。

\begin{table}[H]
\centering
\caption{Replica 各场景标注结果统计}
\label{tab:scene_results}
\begin{tabular}{lccccc}
\toprule
\textbf{场景} & \textbf{LLaVA 原始} & \textbf{过滤保留} & \textbf{精炼描述} & \textbf{空间关系} & \textbf{QA 总数} \\
\midrule
room0 & 74 & 11 & 69 & 32 & 102 \\
room1 & 67 & 48 & 67 & 33 & 101 \\
office0 & 78 & 48 & 23 & 4 & 28 \\
office1 & — & — & 14 & 4 & 19 \\
office2 & — & — & 30 & 10 & 41 \\
office3 & — & — & 34 & 11 & 46 \\
office4 & — & — & 22 & 10 & 33 \\
\midrule
\textbf{合计} & — & — & \textbf{259} & \textbf{104} & \textbf{370} \\
\bottomrule
\end{tabular}
\end{table}

从上表可以观察到以下模式：

\textbf{新旧 prompt 的显著差异。} room0 和 room1 使用优化前的 prompt 生成，QA 数量较高（102 和 101 条），但包含了大量语义空洞的无效应答——如"The perspective gives a sense of depth to the scene"或"An object is the focus of the image"等无法关联到具体物体的描述。而使用优化后 prompt 的 office0—4 场景，虽然绝对 QA 数量较低（平均 33.4 条），但每条 QA 均包含可辨识的物体类别信息，实际可用性显著提高。

\textbf{办公室场景间的较大方差。} office1 仅产出 19 条 QA，而 office3 产出 46 条 QA，两者相差 2.4 倍。这一差异主要源于场景本身的复杂度差异——office1 是一个相对空旷的个人办公室，物体数量和种类较少；office3 则包含更多家具和装饰物，LLaVA 有更多可描述的视觉元素。

\textbf{空间关系数量的衰减。} 优化后 prompt 的 office 场景空间关系平均仅 7.8 条，远低于优化前的 room0（32 条）和 room1（33 条）。这是因为优化后的 prompt 要求更高的关系判定置信度（"Only include pairs with clear spatial relationship"），有效抑制了 LLM 的无根据猜测，但也减少了关系的绝对数量。

\subsection{描述优化前后对比}

表 \ref{tab:caption_quality} 呈现了 prompt 优化前后的典型描述对比，直观展示了新旧 prompt 在语义上的本质差异。

\begin{table}[H]
\centering
\caption{描述优化前后典型示例对比}
\label{tab:caption_quality}
\small
\begin{tabular}{p{6.5cm}p{6.5cm}}
\toprule
\textbf{优化前（旧 prompt）} & \textbf{优化后（新 prompt）} \\
\midrule
The perspective gives a sense of depth to the scene. & (door) wooden, easy to open, against the wall \\
\rule{0pt}{13pt} A floor extends to the edges of the room. & (floor) clean, smooth and well-maintained \\
\rule{0pt}{13pt} Natural light floods the space, creating a bright atmosphere. & (window) glass window allowing light \\
\rule{0pt}{13pt} An object is the focus of the image and illuminated by a light source. & (table) wooden desk with chair \\
\rule{0pt}{13pt} The area is clean and free of any clutter or decorations. & (sofa) couch in living or dining area \\
\bottomrule
\end{tabular}
\end{table}

优化前后的对比揭示了一个质的变化：旧 prompt 的输出全部是"场景级感官描述"——描述的是全局的光照、深度、氛围、清洁度等无法绑定到特定物体的属性；而新 prompt 的输出转变为"物体级实体描述"——每条都包含可辨识的物体类别（door、floor、window、table、sofa），并提供了材质（wooden、glass、fuzzy）、形态（smooth、rectangular）和位置（against the wall）等可直接用于标注的结构化信息。

量化评估表明，新 prompt 将物体类别识别率从接近 0（旧 prompt 输出中几乎没有可辨别的物体类别）提升至约 85\%（新 prompt 输出中的大部分描述可明确关联到具体物体类别）。剩余的约 15\% 包括两类情况：LLaVA 输入图像质量过差导致模型正确标记为"unknown"的，以及模型确实识别错误的少量案例。

\subsection{识别率瓶颈分析}

值得注意的是，在 office0—4 标注的预处理阶段，约 62\% 的 LLaVA 原始描述因截断（如以半截词开头）或语义过短（$<10$ 字符）而被过滤，意味着只有约 38\% 的 LLaVA 输出进入了后续的优化阶段。这一数据直接指向当前管线的首要瓶颈——\textbf{LLaVA 的物体裁剪质量决定了整个管线的标注产量上限}。改进策略包括基于深度信息的自适应裁剪框、多视角描述融合以及更强视觉编码器的引入，已在第五章中详细讨论。

\section{Prompt 消融实验}

为量化本文结构化 prompt 设计中各组件的独立贡献，在 office0 场景的同一组 10 条 LLaVA 原始描述上进行了消融实验。测试了三种 prompt 条件：条件 A（旧 prompt）——简单的"Clean up each description"，无类别输出要求，无禁用词汇表；条件 B（完整新 prompt）——包含类别强制输出、禁用词汇表和容错处理三个组件；条件 C（消融 prompt）——包含类别强制输出和容错处理，但不包含禁用词汇表。

实验结果如表 \ref{tab:ablation} 所示。

\begin{table}[H]
\centering
\caption{Prompt 消融实验结果}
\label{tab:ablation}
\begin{tabular}{lccccc}
\toprule
\textbf{条件} & \textbf{产出数} & \textbf{耗时} & \textbf{类别识别率} & \textbf{禁用词出现} & \textbf{含 class 字段} \\
\midrule
A: 旧 prompt & 10 & 4.5 s & $\sim$0\% & 8 次 & 否 \\
B: 完整新 prompt & 8 & 5.2 s & $\sim$85\% & 0 次 & 是 \\
C: 消融（无禁用词表） & 8 & 5.0 s & $\sim$75\% & 4 次 & 是 \\
\bottomrule
\end{tabular}
\end{table}

消融实验揭示了三项关键发现：

\textbf{类别强制输出是最大贡献组件。} 对比条件 A（旧 prompt，0\% 类别识别率）与条件 C（无禁用词表，75\% 类别识别率），仅增加类别输出要求就将类别识别率从零提升到可用水平。这验证了 Cheng 等人\cite{chain_of_thought}的核心论点——明确的格式约束（强制输出 class 字段）对强大语言模型比提供示例更有效。

\textbf{禁用词汇表提供了额外 10 个百分点的质量增益。} 条件 B 的类别识别率（85\%）高于条件 C（75\%），且禁用词出现次数从 4 次降至 0 次。直接列出禁止使用的词汇（light、atmosphere、depth 等）是最简单但最有效的输出质量控制手段，其效果优于仅靠模型自身的判断。

\textbf{三个组件之间存在协同效应。} 类别强制输出定义了输出的"骨架"（class 字段），禁用词汇表清除了最常见的不良填充物（氛围和光照描述），容错处理为模型无法识别的情况提供了退路（标记为 unknown 而非随意编造）。三者协同使整个 prompt 的质量控制效果大于各组件的简单加和。

\section{合成场景验证实验}

在 Isaac Sim 生成的 20 个合成场景上，基于几何真值自动导出标注。表 \ref{tab:isaac_results} 按房型汇总了标注结果。

\begin{table}[H]
\centering
\caption{Isaac Sim 合成场景标注统计}
\label{tab:isaac_results}
\begin{tabular}{lccccc}
\toprule
\textbf{房型} & \textbf{场景数} & \textbf{平均物体数} & \textbf{平均关系数} & \textbf{总 QA} \\
\midrule
起居室 & 5 & 10.0 & 45.0 & 280 \\
卧室 & 5 & 8.0 & 28.0 & 185 \\
厨房 & 5 & 8.0 & 28.0 & 185 \\
书房 & 5 & 7.0 & 21.0 & 145 \\
\midrule
\textbf{合计} & \textbf{20} & \textbf{8.3} & \textbf{30.5} & \textbf{795} \\
\bottomrule
\end{tabular}
\end{table}

Isaac Sim 的标注结果呈现了两个显著特征。第一，空间关系数量的跨房型差异很大：起居室因含 10 个物体而平均产出 45 条关系（物体数与关系的比值为 1:4.5），书房仅 7 个物体却产出 21 条关系（比值 1:3.0），相对密度反而更高——这是因为书房中物体布局更为紧凑（房间面积 14 m$^2$ 对起居室的 30 m$^2$），更多物体对满足"next to"或"near"的距离阈值。

第二，所有标注均为 100\% 准确——物体的位置从 USD 场景图精确查询（精度 0.01 m），空间关系由确定的几何规则计算（不存在模型推断的随机性）。这与 Concept-Graphs 管线形成鲜明对比：后者通过 LLM 推断的关系存在语法错误（如"Object 7 is creates contrast"）和语义空洞（如将"clean environment details"作为关系理由）。在效率方面，Isaac Sim 单场景生成与标注约需 10 秒，为 Concept-Graphs 管线（约 36 分钟）的 216 倍。

\section{推理模型选型实验}

\subsection{基准测试}

为选择最适合标注管线的推理大语言模型，在单张 RTX 4090 GPU 上对 5 款国内开源模型\cite{qwen2_5}\cite{deepseek_r1}\cite{glm4}进行了横向基准测试。测试任务包括 5 条描述优化和 2 条关系推理，温度参数设为 0.3。表 \ref{tab:model_benchmark} 呈现了测试结果。

\begin{table}[H]
\centering
\caption{5 款大语言模型推理基准测试结果（RTX 4090）}
\label{tab:model_benchmark}
\begin{tabular}{lcccc}
\toprule
\textbf{模型} & \textbf{描述质量} & \textbf{描述速度} & \textbf{关系速度} & \textbf{JSON 合法率} \\
\midrule
Qwen2.5-3B & 100\% & 86.2 t/s & 87.6 t/s & 100\% \\
Qwen2.5-7B & 100\% & 71.0 t/s & 76.0 t/s & 100\% \\
Qwen2.5-14B & 100\% & 45.9 t/s & 51.7 t/s & 100\% \\
DeepSeek-R1-7B & 100\% & 94.2 t/s & 77.5 t/s & 0\% \\
GLM-4-9B & 90\% & 68.0 t/s & 67.5 t/s & 100\% \\
\bottomrule
\end{tabular}
\end{table}

结果分析如下：

\textbf{Qwen2.5-3B 是最优选择。} 在 100\% 语义质量和 JSON 格式合法率的前提下，3B 版本的推理速度（86.2 t/s）比 7B 快 21\%，比 14B 快 88\%。更小的模型参数意味着更低的显存占用（1.9 GB）和更快的加载时间（约 2 秒）。已将其作为管线的默认推理引擎。

\textbf{模型大小的边际收益递减。} 在 Qwen2.5 系列中，描述速度随模型大小呈近似线性下降（3B $\rightarrow$ 7B: $-18\%$，7B $\rightarrow$ 14B: $-35\%$），但语义质量并未随之提升——三个变体的质量评分均为 100\%。这说明在结构化数据合成这类"格式约束强、语义复杂性有限"的任务上，更大的模型并不带来质量增益，反而付出了显著的速度代价。

\textbf{DeepSeek-R1 的思维链机制与结构化输出不兼容。} DeepSeek-R1 在生成正式内容前会自动插入推理过程（\texttt{<think>...</think>} 标签段），虽然这在数学推理等开放式任务中是优势（显式的中间推理有助于验证和纠错），但在结构化 JSON 输出任务中直接导致格式解析完全失败。这一结果提示了一个重要的模型选择原则：\textbf{推理能力和格式遵守在模型行为中可能存在 trade-off}——强化了推理链生成的训练可能会削弱模型对严格输出格式的遵守程度。

\subsection{实际管线对比验证}

基准测试的结果需要在实际管线运行中得到验证。为此，在 office0 场景上分别使用 Qwen2.5-3B 和 Qwen2.5-7B 运行了完整的 Step 5 流程（描述优化 + 关系推理 + QA 生成）。结果如表 \ref{tab:model_pipeline_compare} 所示。

\begin{table}[H]
\centering
\caption{不同模型在实际管线上的产出对比（office0）}
\label{tab:model_pipeline_compare}
\begin{tabular}{lcccc}
\toprule
\textbf{模型} & \textbf{精炼描述} & \textbf{空间关系} & \textbf{QA 总数} & \textbf{用时} \\
\midrule
Qwen2.5-3B & 21 & 4 & 26 & 19 s \\
Qwen2.5-7B & 18 & 3 & 22 & 18 s \\
\bottomrule
\end{tabular}
\end{table}

在实际管线中，3B 模型的描述精炼产出（21 条）略高于 7B 模型（18 条），验证了基准测试的定量结论——在结构化数据合成这类格式约束明确的任务上，更小的模型即可提供足够质量，更大的模型不仅不带来质量增益，反而在速度为王的生产场景中处于劣势。两个模型的实际耗时接近（19 s vs 18 s），因为管线的主要耗时在模型加载和网络通信上，而非推理本身。

\section{管线效率分析}

表 \ref{tab:pipeline_time} 对比了 GPU 加速前后各步骤的实际耗时。端到端管线的总加速比为 1.4 倍，但这一数字有明显的误导性——它被不涉及 LLM 推理的步骤（Step 2 和 Step 3 合计占总时间的 91\%）所稀释。LLM 推理本身的加速比为 18 倍，这对于管线的迭代调试具有显著实用价值：每次调整 prompt 或参数后，验证周期从 15 分钟缩短至 50 秒。

\section{三源数据对比}

表 \ref{tab:three_source} 从多个维度系统对比了三类数据源的特征。

\textbf{互补性分析。} 三源数据呈现出功能互补而非简单替代的关系。Code2Logic 独有的 Strategy Optimization 类型 QA（267 条，涉及多步规划）是静态三维场景无法提供的；Isaac Sim 凭借几何规则产生大量的高精度空间关系（610 条），弥补了 Concept-Graphs 在关系推理准确率上的不足；Concept-Graphs 则提供了真实传感器噪声下的推断基准，是连接纯合成世界与真实物理世界的桥梁。这种互补关系为论文的三层架构提供了直接的实验支撑：每一层解决了一个其他层无法独立解决的问题。

\textbf{效率与精度的权衡。} 三源数据清楚地展现了自动化数据合成中的一个基本权衡：更高的标注精度通常意味着更慢的数据生成速度（Isaac Sim 的几何计算最快但受限于预定义家具库）或更窄的领域覆盖（Code2Logic 的游戏逻辑最精确但限于二维游戏）。Concept-Graphs 管线处于中间位置——通过开放词汇视觉模型实现任意物体描述，但以牺牲精度为代价。这种权衡的量化表征为不同应用场景下的方法选择提供了参考依据。


\chapter{总结与展望}
% 第5章 总结与展望

\section{工作总结}

本文围绕"把游戏代码驱动的数据合成方法迁移到真实 3D 场景"这条主线，在方法迁移、管线构建、推理优化和评估框架四个层面开展了系统性工作。

在方法迁移层面，本文首先验证了 Code2Logic 在 2D 游戏领域的有效性（30 款游戏、1397 条 QA），然后将这套模板化数据合成的思路搬到了三维场景理解任务中。通过 Concept-Graphs 管线，在 Replica 数据集的 7 个室内场景上实现了从 RGB-D 图像序列到结构化 QA 的自动标注。进一步地，为了给推断标注提供可量化的精度基准，搭建了基于 Isaac Sim 的合成真值验证系统，生成了 20 个场景、795 条带有精确 ground truth 的 QA 数据。三层递进——方法验证、真实迁移、真值校准——三类数据累计 2395 条 QA，覆盖五种推理类型。需要指出的是，Code2Logic 的 1397 条 QA 直接使用了 Tong 等人公开发布的 GameQA 数据集，本文的创造性贡献在于自行构建了 Concept-Graphs 管线的 370 条 QA 和 Isaac Sim 的 795 条 QA。

在管线构建层面，完成了一条从 RGB-D 图像序列和相机轨迹到结构化 QA 数据集的六步端到端流水线。输入经过 GradSLAM 三维重建、Grounded-SAM 二维分割、CF-SLAM 三维建图后得到物体级别的场景表示，再由 LLaVA 为每个物体生成自然语言描述，最后由 llama.cpp 在 GPU 上完成描述优化和空间关系推理。管线的关键改进集中在第五步的 prompt 设计：通过强制输出物体类别字段、设置禁用词汇表和引入容错处理，将物体类别识别率从接近零提升至约 85\%。消融实验量化了三个 prompt 组件的独立贡献，其中类别强制输出贡献最大（从 0\% 到 75\%），禁用词汇表额外贡献了约 10 个百分点。同场景新旧 prompt 的对比进一步验证了这一改进的实质效果：虽然新 prompt 的 QA 数量约降到了旧 prompt 的 30\%--35\%，但语义内容发生了从"场景氛围描述"到"物体级标注"的质的改变。

在推理优化层面，排查并解决了 Ollama 内置 CUDA 12.8 库与服务器驱动 565.57 的版本不兼容问题。通过源码编译 llama-cpp-python（链接系统 CUDA 12.6 库），将 qwen2.5 模型的全部层加载至 GPU 显存，推理吞吐量从 3 tok/s 提升至 99 tok/s（33 倍），单场景 LLM 处理时间从 15 分钟降至 50 秒。在 5 款国产开源模型上的横向对比中，Qwen2.5-3B 在速度、质量和格式合法性三个维度上综合最优。一个意外但重要的发现是，14B 模型的实际产出（8 条描述、0 条关系）远低于 3B（21 条描述、4 条关系），说明在结构化输出任务上，更大的模型安全对齐更强，反而导致更多的输出抑制。

在评估框架层面，建立了从数据量、推理类型、关系类型、标注精度和管线效率五个维度系统对比三类数据源的统一评估体系，同时也为其他数据合成方法的比较提供了一个可复用的分析模板。

\section{局限性}

这项工作存在几项需要诚实面对的局限。

最突出的瓶颈来自 LLaVA 的输入质量。7 个场景 539 条 caption 的过滤统计显示，62\% 的原始描述在进入 LLM 优化前就被预处理筛掉了——截断（以半截词开头）占 39\%，氛围词描述占 13\%。无论下游的 prompt 和推理做到什么程度，标注产量的硬上限仅有输入量的约 38\%。这个瓶颈的根源在视觉语言模型的物体裁剪阶段：当裁剪区域太小或背景噪声过多时，LLaVA 无法生成有意义的物体描述。

此外，场景多样性有限制了泛化性评估的充分性。当前的 7 个 Replica 室内场景（两个起居室、五个办公室）和 Isaac Sim 的 4 种房型模板全部属于室内静态环境，不涉及室外、工业或含有动态物体的场景。

还有两个需要后续补全的验证环节。第一，本文的实验集中于标注数据的产出过程，没有将 2395 条 QA 用于下游 VLM 的训练并评估实际效果。标注数据对模型性能的贡献是检验数据质量的终极标准，这个实验闭环在本研究中尚未完成。第二，Isaac Sim 的渲染和 ground truth 导出验证通过了，CG 管线的各步骤也在 7 个场景上独立验证了稳定性，但两个系统之间的完整自动串联流程还没有跑通。

\section{未来方向}

针对上述局限，后续工作有几个明确的方向。

在输入质量方面，基于深度信息的自适应裁剪可以在物体远近不同时动态调整裁剪框大小，兼顾近距离物体的上下文保留和远距离物体的噪声抑制。多视角描述融合也能有效滤除单帧噪声——同一物体在多个观测帧中生成的独立描述，经过交叉验证后的一致性筛选可以大幅降低幻觉率。An 等人\cite{multimodal_3d} 提出的深度与 RGB 特征联合编码方法是可直接借鉴的技术路线。

在下游验证方面，将 2395 条 QA 用于微调 Qwen2.5-VL 等视觉语言模型，在 MathVista、MMMU、GQA 等标准基准上评估微调前后的性能变化，是检验数据质量最直接的方式。Code2Logic 原始实验中使用的 GRPO 强化学习策略\cite{code2logic}——分别降低了 19.2\% 的感知错误率和 11.5\% 的推理错误率——可以引入本文的数据训练流程。

在场景扩展方面，ScanNet 和 HM3D 等更大规模的数据集是真实场景覆盖面的直接补充。Isaac Sim 中可以构建更多样化的房型模板，引入动态物体和时序连续渲染，评估管线在变化环境下的性能。Song 等人\cite{llm_planner} 关于具身代理在动态环境中多步规划挑战的研究，为动态场景的评估范式提供了参考。

真值验证链路的全自动化同样值得继续推进。将 Isaac Sim 渲染、CG 管线标注和真值对比三个环节串联后，可以在物体检测数量、类别准确率、三维位置误差和关系准确率四个指标上直接给出精确的量化结果，使"标注精度可验证"这一主张从概念落地为可操作的评估流程。


\backmatter
\printbibliography

\begin{acknowledgements}
本论文的完成离不开许多人的帮助与支持。

首先，衷心感谢我的指导教师戈维峰副教授。在整个毕业设计过程中，戈老师从选题方向、研究方法到论文撰写规范都给予了悉心指导和宝贵建议。

感谢计算与智能创新学院的各位老师，在四年本科学习期间传授的专业知识和科研方法，为本论文的研究工作奠定了坚实的基础。

感谢开源社区的贡献者。本文的管线构建依赖于 GradSLAM、Grounded-SAM、LLaVA、llama.cpp、Ollama 等优秀的开源项目，Code2Logic 论文作者公开的 GameQA 数据集为本研究提供了重要的方法论参考和对比基线。NVIDIA Isaac Sim 团队提供了功能完备的仿真平台。

感谢我的同学和朋友们，在毕设过程中与我讨论技术问题、分享经验心得，在遇到困难时给予鼓励与支持。

最后，感谢我的家人多年来的关怀与付出，你们的理解与支持是我不断前行的动力。
\end{acknowledgements}

\appendix

\chapter{完整 Prompt 模板}

以下为本文在描述优化和关系推理阶段使用的完整 prompt 文本。

\section*{A.1 描述优化 System Prompt}

\begin{verbatim}
You analyze object crops from an indoor 3D scene.
For each item, identify the object and output a concise description.

Output ONLY a JSON array:
[{"id": <number>, "class": "<category>", "caption": "<description>"}]

RULES for caption:
- Name the object category (door, floor, ceiling, window, table,
  chair, sofa, cabinet, shelf, light, curtain, pillow, rug, etc.)
- Describe material/color/shape (e.g. "wooden", "white", "rectangular")
- State its position if evident (e.g. "against the wall", "in the corner")
- Keep it to ONE sentence, under 20 words

FORBIDDEN - never use these words:
  light, lighting, shadow, bright, dark, depth, atmosphere,
  airy, mood, perspective, contrast, illuminate

If the input text is too vague to identify the object,
set class="unknown" and caption="An unidentifiable object."

If the text is clearly truncated/incomplete, skip the item entirely.
\end{verbatim}

\section*{A.2 关系推理 Prompt}

\begin{verbatim}
You are given a list of objects with their positions in a room.
Determine spatial relationships between pairs that are clearly related.

Output ONLY a JSON array:
[{"from": <id1>, "to": <id2>, "relation": "<type>", "reason": "<brief>"}]

Relation types: "next to" | "above" | "below" | "behind" |
                "in front of" | "on"

Only include pairs that have a clear spatial relationship. Do not guess.
\end{verbatim}

\chapter{QA 数据 JSON Schema}

\section*{B.1 Target Perception（目标感知）}

\begin{verbatim}
{
  "data_id": "<scene>-tp-<id:04d>",
  "qa_type": "Target Perception",
  "question": "What is object <id> in this scene?",
  "answer": "<refined caption>"
}
\end{verbatim}

\section*{B.2 Spatial Reasoning（空间推理）}

\begin{verbatim}
{
  "data_id": "<scene>-sp-<src:04d>-<dst:04d>",
  "qa_type": "Spatial Reasoning",
  "question": "Where is object <dst> relative to object <src>?",
  "answer": "Object <dst> is <relation> object <src>."
}
\end{verbatim}

\section*{B.3 Scene Comprehension（场景理解）}

\begin{verbatim}
{
  "data_id": "<scene>-sc",
  "qa_type": "Scene Comprehension",
  "question": "Describe the overall layout of this scene.",
  "answer": "A <room_type> with <N> objects: <summary>..."
}
\end{verbatim}

\end{document}
